% BHU PYQ -- Transcribed & Corrected
% Paper: SANMJ-31 / SANMJ31: Sanskrit Sahitya-II (Mahakavya)
% Exam: B.A. (Hons.) Semester III Examination, 2025-26 (NEP)
% Category: Major Core
% Department: Department of Sanskrit, Faculty of Arts, Banaras Hindu University

\documentclass[11pt,a4paper]{article}
\usepackage[a4paper,top=2.5cm,bottom=2.5cm,left=2.8cm,right=2.8cm,headheight=14pt]{geometry}
\usepackage{amsmath,amssymb}
\usepackage[T1]{fontenc}
\usepackage[utf8]{inputenc}
\usepackage{lmodern,microtype,fancyhdr,enumitem,hyperref,lastpage,parskip}

\hypersetup{
    colorlinks=true,
    urlcolor=black,
    linkcolor=black,
    pdftitle={SANMJ31: Sanskrit Sahitya II (Mahakavya) -- BHU B.A. Sem III 2025-26 (NEP)}
}

\pagestyle{fancy}
\fancyhf{}
\renewcommand{\headrulewidth}{0.4pt}
\renewcommand{\footrulewidth}{0.4pt}
\fancyhead[L]{\small \textit{Banaras Hindu University}}
\fancyhead[R]{\small \textit{SANMJ-31: Sanskrit Sahitya II (Sem III 2025-26)}}
\fancyfoot[C]{\small Page \thepage\ of \pageref{LastPage}}

\newlist{parts}{enumerate}{2}
\setlist[parts,1]{label=(\alph*),leftmargin=*,itemsep=4pt,topsep=2pt}

\newcommand{\pts}[1]{\hfill \textit{[#1]}}

\begin{document}

\begin{center}
  {\large\bfseries BANARAS HINDU UNIVERSITY}\\[2pt]
  {\normalsize B.A. (Hons.) Semester III Examination, 2025-26 (NEP)}\\[6pt]
  \rule{0.8\linewidth}{0.5pt}\\[6pt]
  {\large\bfseries DEPARTMENT OF SANSKRIT}\\[4pt]
  {\large\bfseries Paper Code: SANMJ-31 : संस्कृत साहित्य-II (महाकाव्य)}\\[2pt]
  {\small Ref. Code: 2025-26/D-III/055}\\[4pt]
  {\normalsize Time: 2:30 Hours \hfill Full Marks: 70}
\end{center}

\rule{\linewidth}{0.4pt}

\smallskip

\noindent \textit{(Write your Roll No. at the top immediately on the receipt of this question paper)}\\[2pt]
\noindent \textit{सूचना : सभी खण्डों के प्रश्नों के उत्तर निर्देशानुसार दीजिए।}

\bigskip

\begin{center}
  {\large\bfseries SECTION A / खण्ड-अ}\\[2pt]
  {\normalsize (Long Answer Type Questions / दीर्घ उत्तरीय प्रश्न)}
\end{center}

\noindent \textit{सूचना : किन्हीं तीन प्रश्नों के उत्तर दीजिए। प्रत्येक प्रश्न 10 अंकों का है।}\\[2pt]

\bigskip

\noindent \textbf{Question 1.} महाकाव्य के लक्षणों का सोदाहरण वर्णन कीजिए। \pts{10}

\medskip\rule{\linewidth}{0.2pt}\medskip

\noindent \textbf{Question 2.} 'किरातार्जुनीयम्' के प्रथम सर्ग के आधार पर वनेचर की कथन-शैली एवं चरित्र की विशेषताओं की विवेचना कीजिए। \pts{10}

\medskip\rule{\linewidth}{0.2pt}\medskip

\noindent \textbf{Question 3.} 'शिशुपालवधम्' के प्रथम सर्ग के आधार पर श्रीकृष्ण एवं नारद संवाद की सोदाहरण समीक्षा कीजिए। \pts{10}

\medskip\rule{\linewidth}{0.2pt}\medskip

\noindent \textbf{Question 4.} महाकवि भारवि की अर्थगौरव-प्रियता पर विस्तृत लेख लिखिए। \pts{10}

\medskip\rule{\linewidth}{0.2pt}\medskip

\noindent \textbf{Question 5.} 'शिशुपालवधम्' के प्रथम सर्ग में वर्णित नारद के व्यक्तित्व एवं रूप-सौन्दर्य का वर्णन कीजिए। \pts{10}

\bigskip
\rule{\linewidth}{0.2pt}
\bigskip

\begin{center}
  {\large\bfseries SECTION B / खण्ड-ब}\\[2pt]
  {\normalsize (Short Answer Type Questions / लघु उत्तरीय प्रश्न)}
\end{center}

\noindent \textit{सूचना : किन्हीं चार प्रश्नों के उत्तर दीजिए। प्रत्येक प्रश्न 5 अंकों का है।}\\[2pt]

\bigskip

\noindent \textbf{Question 6.} भारवेरर्थगौरवम् पर संक्षिप्त टिप्पणी लिखिए। \pts{5}

\medskip\rule{\linewidth}{0.2pt}\medskip

\noindent \textbf{Question 7.} वनेचर द्वारा वर्णित दुर्योधन की शासन-पद्धति पर टिप्पणी लिखिए। \pts{5}

\medskip\rule{\linewidth}{0.2pt}\medskip

\noindent \textbf{Question 8.} 'नवसर्गगते माघे नवशब्दो न विद्यते' कथन की समीक्षा कीजिए। \pts{5}

\medskip\rule{\linewidth}{0.2pt}\medskip

\noindent \textbf{Question 9.} किरातार्जुनीयम् के प्रथम सर्ग में प्रयुक्त छन्दों की सोदाहरण विवेचना कीजिए। \pts{5}

\medskip\rule{\linewidth}{0.2pt}\medskip

\noindent \textbf{Question 10.} शिशुपालवधम् के प्रथम सर्ग के कथासार का वर्णन कीजिए। \pts{5}

\medskip\rule{\linewidth}{0.2pt}\medskip

\noindent \textbf{Question 11.} ससन्दर्भ व्याख्या कीजिए:\\[2pt]
\begin{center}
\textit{``अथावतारोन्मुखमुत्तरीयकं मुनेर्वपुर्विसारिभि: शौचमिवाथ लम्भयन्‌।}\\
\textit{वितानमुच्चैर्व्यतिषज्जयन्निव प्रियं जगद्गद्गदमुच्चैर्वाचमवोचदच्युत:॥''}
\end{center} \pts{5}

\bigskip
\rule{\linewidth}{0.2pt}
\bigskip

\begin{center}
  {\large\bfseries SECTION C / खण्ड-स}\\[2pt]
  {\normalsize (Very Short Answer Type Questions / अतिलघु उत्तरीय प्रश्न)}
\end{center}

\noindent \textit{सूचना : सभी दस भागों के उत्तर दीजिए। प्रत्येक भाग 2 अंकों का है।}\\[2pt]

\bigskip

\noindent \textbf{Question 12.}
\begin{parts}
  \item महाकवि भारवि का समय क्या है? \pts{2}
  \item 'किरातार्जुनीयम्' में कुल कितने सर्ग हैं? \pts{2}
  \item 'किरातार्जुनीयम्' का उपजीव्य ग्रन्थ कौन-सा है? \pts{2}
  \item महाकवि माघ के पिता का क्या नाम था? \pts{2}
  \item 'शिशुपालवधम्' में किस रस की प्रधानता है? \pts{2}
  \item बृहतत्रयी के अन्तर्गत आने वाले तीनों महाकाव्यों के नाम लिखिए। \pts{2}
  \item 'श्रिय: कुरुणामधिपस्य पालनीम्' श्लोक में कौन-सा छन्द है? \pts{2}
  \item माघ के अनुसार काव्य में कौन-से तीन गुण होने चाहिए? \pts{2}
  \item 'वनेचर' शब्द में कौन-सा समास है? \pts{2}
  \item 'शिशुपालवधम्' का मंगलाचरण किस प्रकार का है? \pts{2}
\end{parts}

\vfill
\rule{\linewidth}{0.4pt}

\begin{center}\small
\href{https://bhu-exam.vercel.app/}{Click here to visit our website}\\[4pt]
\href{https://me-aryan.vercel.app/}{Click here to visit our contributor}
\end{center}

\end{document}
